\documentclass[aspectratio=169, 10pt]{beamer}

\usepackage{bm} % bold math
\usepackage{fontspec}
\usepackage{minted}
\usepackage{pgf-pie}
\usepackage{tikz}
\usepackage{graphicx}
\newcommand\sbullet[1][.5]{\mathbin{\vcenter{\hbox{\scalebox{#1}{$\bullet$}}}}}

% Custom commands and environments
\makeatletter
\newcommand\version[1]{\renewcommand\@version{#1}}
\newcommand\@version{}
\def\insertversion{\@version}

\newcommand\course[1]{\renewcommand\@course{#1}}
\newcommand\@course{}
\def\insertcourse{\@course}

\newcommand\coursetitle[1]{\renewcommand\@coursetitle{#1}}
\newcommand\@coursetitle{}
\def\insertcoursetitle{\@coursetitle}

\newcommand\lecturenumber[1]{\renewcommand\@lecturenumber{#1}}
\newcommand\@lecturenumber{}
\def\insertlecturenumber{\@lecturenumber}
\makeatother

\newcommand{\slidetitle}[1]{{\xbseries \large \structure{#1}} \bigskip}
\newcommand{\term}[1]{{\color{blue} #1}}
\newcommand{\leftspace}{\hspace{1em}}
\newcommand{\inlinearrow}{
  \begin{tikzpicture}[baseline]
    \node [anchor=base] (x) {};
    \draw [rawarrow] (x.mid west) -- ($(x.mid west) + (2em,0)$);
  \end{tikzpicture}
}

\newenvironment{slide}
{\begin{frame}[fragile,environment=slide]\vskip0pt plus 1filll}
{\vskip0pt plus 1filll\end{frame}}

% LaTeX

\setlength{\leftmargini}{1em}

% Common Information

\author{Talia Xu}
\course{COMPSCI 340}
\coursetitle{Operating Systems}
\date{2024 Semester 2}

% fontspec

\defaultfontfeatures{Ligatures=TeX}
% \setmainfont{Domine}
\setsansfont{Inter}[
  FontFace={ul}{n}{Font=*-Thin},
  FontFace={el}{n}{Font=*-ExtraLight},
  FontFace={l}{n}{Font=*-Light},
  FontFace={sb}{n}{Font=*-SemiBold},
  FontFace={eb}{n}{Font=*-ExtraBold},
  FontFace={xb}{n}{Font=*-Black},
]
\setmonofont[Contextuals=AlternateOff, Ligatures=TeXOff]{Iosevka}[
  FontFace={xb}{n}{Font=*-Heavy},
]

%% Font Weights

\DeclareRobustCommand{\ulseries}{\fontseries{ul}\selectfont}
\DeclareTextFontCommand{\textul}{\ulseries}
\DeclareRobustCommand{\elseries}{\fontseries{el}\selectfont}
\DeclareTextFontCommand{\textel}{\elseries}
\DeclareRobustCommand{\lseries}{\fontseries{l}\selectfont}
\DeclareTextFontCommand{\textl}{\lseries}
\DeclareRobustCommand{\sbseries}{\fontseries{sb}\selectfont}
\DeclareTextFontCommand{\textsb}{\sbseries}
\DeclareRobustCommand{\ebseries}{\fontseries{eb}\selectfont}
\DeclareTextFontCommand{\texteb}{\ebseries}
\DeclareRobustCommand{\xbseries}{\fontseries{xb}\selectfont}
\DeclareTextFontCommand{\textxb}{\xbseries}

% tikz

\usetikzlibrary{
  arrows,
  arrows.meta,
  automata,
  backgrounds,
  calc,
  decorations.pathreplacing,
  matrix,
  positioning,
  overlay-beamer-styles,
  shapes,
  shapes.multipart,
  tikzmark,
}

\tikzstyle{rawarrow} = [
  -{Latex[round]},
  line width=1pt,
  yellow,
  shorten >=3pt,
  shorten <=3pt,
  font=\small,
  text=black,
]

\tikzstyle{arrow} = [
  -{Latex[round]},
  line width=1pt,
  yellow,
  shorten >=3pt,
  shorten <=3pt,
  transform canvas={yshift=3pt},
  font=\small,
  text=black,
]

\newcommand{\tikzmarkcoord}[1]{([yshift=3pt]pic cs:#1)}

% minted

\setminted{style=eyolfson, fontsize=\small, escapeinside=||}
\setmintedinline{fontsize=\normalsize}

% hyperref

\hypersetup{colorlinks, urlcolor=blue}

% beamer
\setbeamersize{text margin left=16mm, text margin right=16mm}
\setbeamertemplate{itemize items}[circle]
\setbeamercolor{item}{fg=black}
\setbeamercolor{structure}{fg=darkblue}
\setbeamerfont{frametitle}{series=\bfseries, parent=structure}
\setbeamertemplate{navigation symbols}{}
\setbeamertemplate{headline}{}
\setbeamertemplate{footline}{
  \begin{tikzpicture}[
    remember picture,
    overlay,
    shift={(current page.south west)},
  ]
    \path [fill=gray] (144mm, 0) -- (160mm, 16mm) -- (160mm, 0);
    \node [inner sep=3.5mm, outer sep=0, text=black, anchor=base east,
           align=right, yshift=3.5mm]
          at (current page.south east) {\ttfamily \small \insertframenumber{}};
  \end{tikzpicture}
}
\setbeamertemplate{title page}{
  \begin{tikzpicture}[
    remember picture,
    overlay,
    shift={(current page.south west)},
    background rectangle/.style={fill=darkblue},
    show background rectangle,
  ]
    \node [anchor=center, align=center, text=white, text width=40mm, scale=3.2]
          at (\paperwidth / 2, \paperheight * 2 / 3)
          {\xbseries \inserttitle{}};
    \node [anchor=base west, align=left, inner sep=0, text=white, yshift=2.5mm]
          at (16mm, \paperheight / 3)
          {\insertdate{} \insertcourse{}: \insertcoursetitle{}};
    \node [anchor=base west, align=left, inner sep=0, text=white, yshift=-2.5mm]
          at (16mm, \paperheight / 3)
          {\insertauthor};
    \node [anchor=base east, align=right, inner sep=0, text=white, yshift=2.5mm]
          at (144mm, \paperheight / 3)
          {Lecture \insertlecturenumber{}};
    \node [anchor=base east, align=right, inner sep=0, text=white,
           yshift=-2.5mm]
          at (144mm, \paperheight / 3)
          {\ttfamily \insertversion{}};
    \node [align=center, anchor=south, inner sep=0, text=white, yshift=3.5mm]
          (license) at (\paperwidth / 2, 0)
          {\fontsize{7pt}{7pt}\selectfont This  work is licensed under a
           \href{http://creativecommons.org/licenses/by-sa/4.0/}
                {\color{lightblue} Creative Commons Attribution-ShareAlike 4.0
                 International License}};
  \end{tikzpicture}
}

% xcolor

%% Primary Colour

\definecolor{pantone655}{RGB}{0, 42, 92} % #002a5c
\colorlet{darkblue}{pantone655}

%% Secondary Colours

\definecolor{pantone633}{RGB}{0, 139, 176} % #008bb0
\colorlet{blue}{pantone633}

\definecolor{pantonewarmred}{RGB}{220, 70, 51} % #dc4633
\colorlet{red}{pantonewarmred}

\definecolor{pantone3285}{RGB}{0, 161, 137} % #00a189
\colorlet{cyan}{pantone3285}

\definecolor{pantone7722}{RGB}{13, 83, 77} % #0d534d
\colorlet{darkcyan}{pantone7722}

\definecolor{pantone376}{RGB}{141, 191, 46} % #8dbf2e
\colorlet{green}{pantone376}

\definecolor{pantone2613}{RGB}{109, 36, 122} % #6d247a
\colorlet{violet}{pantone2613}

\definecolor{pantone2985}{RGB}{111, 199, 234} % #6fc7ea
\colorlet{lightblue}{pantone2985}

\definecolor{pantone227}{RGB}{171, 19, 104} % #ab1368
\colorlet{magenta}{pantone227}

\definecolor{pantone7406}{RGB}{241, 197, 0} % #f1c500
\colorlet{yellow}{pantone7406}

%% Neutrals

\definecolor{pantonecoolgray2}{RGB}{208, 209, 201} % #d0d1c9
\colorlet{gray}{pantonecoolgray2}

\usepackage{svg}

\lecturenumber{26}
\title{Disks}
\version{2.0.0}

\begin{document}

\begin{frame}[plain, noframenumbering]
  \titlepage
\end{frame}

\begin{slide}
  
  \slidetitle{Solid State Drives (SSD) Are More Modern}

  Use transistors (like RAM) to store data rather than magnetic disks
  \medskip

  Pros

  \leftspace{}No moving parts or physical limitations

  \leftspace{}Higher throughput, and good random access

  \leftspace{}More energy efficient

  \leftspace{}Better space density
  \medskip

  Cons

  \leftspace{}More expensive

  \leftspace{}Lower endurance (number of writes)

  \leftspace{}More complicated to write drivers for

\end{slide}

\begin{slide}
  
  \slidetitle{A SSD Contains Pages}

  \begin{center}
    \includegraphics[height=0.8\textheight]{ssd.png}    
  \end{center}

\end{slide}

\begin{slide}
  
  \slidetitle{SSDs Using NAND Flash Are Much Faster Than HDDs}

  Pages are typically 4 KiB
  \medskip

  Reading a page: 10 μs

  Writing a page: 100 μs

  Erasing a block: 1 ms
\end{slide}

\begin{slide}
  
  \slidetitle{NAND Flash Programming Uses Pages and Blocks}

  You can only read complete pages and write to freshly erased pages
  \medskip

  Erasing is done per block (a block has 128 or 256 pages)

  \leftspace{}An entire block needs to be erased before writing
  \medskip

  Writing is slow (may need to create a new block)

\end{slide}

\begin{slide}
  
  \slidetitle{The OS Can Help Speed Up SSDs}

  SSDs need to garbage collect blocks

  \leftspace{}Move any pages that are still alive to a new block (may be overhead)
  \medskip

  The disk controller doesn't know what blocks are still alive

  \leftspace{}SSD may think the disk is full, when a file could be deleted (not erased)
  \medskip

  The OS can use the \textit{TRIM} command to inform the SSD a block is unused

  \leftspace{}The SSD can freely erase the block without moving overhead

\end{slide}

\begin{slide}
  
  \slidetitle{So Far We've Been Talking About Single Devices}

  Sometimes called Single Large Expensive Disk (SLED)

  \leftspace{}Just one large disk for data

  \leftspace{}\leftspace{}Single point of failure
  \medskip

  There's also Redundant Array of Independent Disks (RAID)

  \leftspace{}Data distributed on multiple disks

  \leftspace{}\leftspace{}Use redundancy to prevent data loss

  \leftspace{}\leftspace{}Use redundancy to increase throughput

\end{slide}

\begin{slide}
  
  \slidetitle{RAID 0 is Called a Striped Volume}

  Data stripes (128KB and 256KB) are distributed over disks

  \begin{center}
    \includesvg[height=5cm]{raid-0.svg}
  \end{center}
  
  \begin{flushright}
    by Cburnett licensed under CC BY-SA 3.0
  \end{flushright}

\end{slide}

\begin{slide}
  
  \slidetitle{RAID 0 is For Performance Only}

  The data is stripped across all disks in the array (you can have more than 2)
  \medskip

  Pro

  \leftspace{}Faster parallel access, roughly $\mathsf{N}$ times speed
  \medskip

  Con

  \leftspace{}Any disk failure results in a data loss (more points of failure)

\end{slide}

\begin{slide}
  
  \slidetitle{RAID 1 Mirrors All Data Across All Disks}

  \begin{center}
    \includesvg[height=5cm]{raid-1.svg}
  \end{center}

  \begin{flushright}
    by Cburnett licensed under CC BY-SA 3.0
  \end{flushright}
\end{slide}

\begin{slide}
  
  \slidetitle{RAID 1 is Simple, But Wasteful}

  Every disk in the array has a mirrored copy of all the data
  \medskip

  Pro

  \leftspace{}Good reliability, as long as one disk remains, no data loss

  \leftspace{}Good read performance
  \medskip

  Con

  \leftspace{}High cost for redundancy (we can do better)

  \leftspace{}Write performance is the same as a single disk

\end{slide}

\begin{slide}
  
  \slidetitle{RAID 4 Introduces Parity}

  Data stripes distributed over disks with a dedicated parity disk (p = parity)

  \leftspace{}Parity stores xor $\oplus$ of copies 1-3, any one copy can be
               reconstructed

  \begin{center}
    \includesvg[height=5cm]{raid-4.svg}
  \end{center}

  \begin{flushright}
    by Cburnett licensed under CC BY-SA 3.0
  \end{flushright}

\end{slide}

\begin{slide}
  
  \slidetitle{RAID 4 Can Use the Parity Drive to Recover}

  With parity, we can use $\mathsf{1 - \frac{1}{N}}$ of the available space

  \leftspace{}Requires at least 3 drives
  \medskip

  Pro

  \leftspace{}We get $\mathsf{(N - 1)}$ times performance
               (removing parity disk)

  \leftspace{}We can replace a failed disk and rebuild
  \medskip

  Con

  \leftspace{}Write performance can suffer, every write must write to parity
               disk

\end{slide}

\begin{slide}
  
  \slidetitle{RAID 5 Distributes Parity Across All Disks}

  Data stripes distributed over disks and each disk takes turns with parity
  blocks

  \begin{center}
    \includesvg[height=5cm]{raid-5.svg}
  \end{center}

  \begin{flushright}
    by Cburnett licensed under CC BY-SA 3.0
  \end{flushright}

\end{slide}

\begin{slide}
  
  \slidetitle{RAID 5 is an Improved Raid 4}

  It has all the same pros as RAID 4
  \medskip

  Write performance is improved, no longer a bottleneck on a single parity
  drive

\end{slide}

\begin{slide}
  
  \slidetitle{RAID 6 Adds Another Parity Block Per Stripe}

  \begin{center}
    \includesvg[height=5cm]{raid-6.svg}
  \end{center}

  \begin{flushright}
    by Cburnett licensed under CC BY-SA 3.0
  \end{flushright}

\end{slide}

\begin{slide}

  \slidetitle{RAID 6 Can Recover from 2 Simultaneous Drive Failures}

  Due to the extra parity, we can use $\mathsf{1 - \frac{2}{N}}$ of the
  available space

  \leftspace{}Requires at least 4 drives
  \medskip

  Write performance is slightly less than RAID 5, due to another parity
  calculation

\end{slide}

\begin{slide}
  
  \slidetitle{Disks Enable Persistence}

  We explored two topics: SSDs and RAID

  \begin{itemize}
    \item SSDs are more like RAM except accessed in pages and blocks
    \item SSDs also need to work with the OS for best performance (TRIM)
    \item Use RAID to tolerate failures and improve performance using multiple disks
  \end{itemize}

\end{slide}

\end{document}
